% Isolated repair of Steps 2--3 in prop:ch23-model-vs-psh-envelope.
% Preserve the original fixed skeleton; no new determining-model choice.
\emph{Step 2: approximation on the fixed skeleton.}
By clause (2) of \Cref{def:ch21-bfj-psh}, the restriction
\(\psi|_{\Delta_{\mathcal Z}}\) is a uniform limit of restrictions of
\(\theta\)-psh real model functions.  Choose such a function \(\chi\)
with
\begin{equation}\label{eq:ch23-skeletal-approximation-error}
 \bigl|\chi-\psi\bigr|<\varepsilon
 \qquad\text{on }\Delta_{\mathcal Z}.
\end{equation}
Keep \(\mathcal Z\) fixed.  The function \(\chi\) need not be determined
on this model; enlarging the skeleton after choosing the approximation
would not preserve the asserted approximation on the enlarged skeleton.

\emph{Step 3: pushing the inequality off the skeleton.}
Because \(\mathcal Z\) determines \(\theta\),
\Cref{lem:ch21-model-skeletal-properties} gives
\(\chi\leq\chi\circ p_{\mathcal Z}\).  Consequently, for every
\(x'\in X^{\An}\),
\[
 \chi(x')\leq\chi(p_{\mathcal Z}(x'))
 \overset{\eqref{eq:ch23-skeletal-approximation-error}}{<}
 \psi(p_{\mathcal Z}(x'))+\varepsilon
 \leq f(p_{\mathcal Z}(x'))+\varepsilon
 \overset{\eqref{eq:ch23-obstacle-retraction-error}}{<}
 f(x')+2\varepsilon.
\]
Thus \(\chi-2\varepsilon\) is a \(\theta\)-psh real model
competitor below \(f\).  The real-model convention in this chapter
allows the real constant \(2\varepsilon\).  In a rational-model
argument it would instead be necessary to subtract an admissible
model constant slightly larger than \(2\varepsilon\).

